\section{Properties of Logarithms}
The logarithm has many properties that allow us to manipulate it in convenient ways. In fact, historically, these rules came \emph{first}; it only turned out later that the function with these properties is the inverse of an exponential. However, we're very familiar with exponents by the time we get to algebra, so introducing the logarithm as an inverse comes more naturally.

These properties seem like a came at this point, but we make significant use of them in solving equations. They will be much more important in calculus, where they can be used to turn incredibly tedious problems in to incredibly easy ones.

Colloquially, we call these properties ``Log Rules.''

\subsection{Exponent Properties for Logarithms} 
The familiar rules for exponents can be rewritten as a rule for logarithms. In particular, we will be interested in the rules involving powers of the same \makeblank[1in]:
\begin{theorem}[Rules of Exponents]
For any real numbers $b$, $m$, $n$, and $p$ where $b>0$:
\begin{itemize}
\item $b^m\cdot b^n = b^{m+n}$ (product of powers rule)
\item $\dfrac{b^m}{b^n} = b^{m-n}$ (quotient of powers rule)
\item $\left(b^m\right)^p=b^{mp}$ (power of a power rule)
\end{itemize}
\end{theorem}
\begin{proof}
These rules are easy to show for \makeblank exponents. Then we defined all other exponents (negative integers, fractions, real numbers\textsuperscript{*}) to make the rules continue to hold.
\end{proof}
To get the same rules for logarithms, we will define:
\[
x=b^m
\qquad
y=b^n
\]
This means that $\log_b(x)=\makeblanksmall$ and $\log_b(y)=\makeblanksmall$.
\vfill
\textsuperscript{*} You might be thinking we didn't really define irrational exponents... and you're right. We use the fact that we can approximate any irrational number by a terminating decimal---which is rational---as closely as we'd like. 
\subsection{Log of a Product} 
\begin{theorem}[Log of a Product Rule]
For any positive real numbers $x$, $y$, and $b$, $b\neq 1$:
\[
\log_b(xy)=\log_b x+\log_b y
\]
\end{theorem}
\begin{proof}
$x$, $y$, and $b$ are positive real numbers, so $x=b^m$ and $y=b^n$ for some real numbers $m$ and $n$.
\begin{itemize}
\item $\log_b(x)=\makeblanksmall$
\item $\log_b(y)=\makeblanksmall$
\item $xy=\blankparens\blankparens=\makeblank[1in]$
\item So $\log_b(xy)=\makeblank[1in]=\makeblank[1in]+\makeblank[1in]$
\end{itemize}
This proves the theorem.
\end{proof}
How do we use this?
\begin{example}
Rewrite $\log_4(5x^2)$ as separate logs.
\begin{lpic}{}{1}
\end{lpic}
\end{example}
\begin{example}
Rewrite $\log_3(x)+\log_3(y+1)$ as a single log.
\begin{lpic}{}{1}
\end{lpic}
\end{example}
We refer to this rule as allowing us to ``split up'' or ``combine'' the logs. 
\subsection{Log of a Quotient} 
\begin{theorem}[Log of a Quotient Rule]
For any positive real numbers $x$, $y$, and $b$, $b\neq 1$:
\[
\log_b\left(\frac{x}{y}\right)=\log_b x-\log_b y
\]
\end{theorem}
\begin{proof}
$x$, $y$, and $b$ are positive real numbers, so $x=b^m$ and $y=b^n$ for some real numbers $m$ and $n$.
\begin{itemize}
\item $\log_b(x)=\makeblanksmall$
\item $\log_b(y)=\makeblanksmall$
\item $\dfrac{x}{y}=\dfrac{\blankparens}{\blankparens}=\makeblank[1in]$
\item So $\log_b\left(\dfrac{x}{y}\right)=\makeblank[1in]=\makeblank[1in]-\makeblank[1in]$
\end{itemize}
This proves the theorem.
\end{proof}
How do we use this?
\begin{Example}
Rewrite $\log_5\left(\dfrac{28}{q}\right)$ as separate logs.
\begin{lpic}{}{1}
\end{lpic}
\end{Example}
\begin{Example}
Rewrite $\log_7(x+1)-\log_7(y^2)$ as a single log.
\begin{lpic}{}{1}
\end{lpic}
\end{Example}
This rule also lets us ``split up'' or ``combine'' logs. 
\subsection{Log of a Power} 
\begin{theorem}[Log of a Product Rule]
For any positive real numbers $x$ and $b$, $b\neq 1$, and any real number $p$:
\[
\log_b\left(x^p\right)=p \log_b x
\]
\end{theorem}
\begin{proof}
$x$ and $b$ are positive real numbers, so $x=b^m$ for some real number $m$.
\begin{itemize}
\item $\log_b(x)=\makeblanksmall$
\item $x^p=\makeblank[1in]=\makeblank[1in]$
\item So $\log_b\left(x^p\right)=\makeblank[1in]=\makeblank[1in]$
\end{itemize}
This proves the theorem.
\end{proof}
How do we use this?
\begin{Example}
Rewrite $\log_3\left(x^5\right)$ as a simpler log.
\begin{lpic}{}{1}
\end{lpic}
\end{Example}
\begin{Example}
Rewrite $\log_3\left(\sqrt[5]{z}\right)$ as a simpler log.
\begin{lpic}{}{2}
\end{lpic}
\end{Example}
\begin{Example}
Rewrite $3\log(x+1)$ with everything inside the log.
\begin{lpic}{}{1}
\end{lpic}
\end{Example}
We refer to ``pulling the power out of'' or ``pushing the power into'' the log. 
\subsection{Log Rules to Separate} 
One use for the log rules is to rewrite the log of a complex expression as a combination of logs of much simpler expressions.
\begin{note}
In the log rules, we assume that every expression whose log we take is positive, since the log is only defined for positive numbers. Sometimes, we will include an absolute value to make this less of a concern. In our work with log rules in this class, we will assume the values of all variables are such that there is no concern about domain.
\end{note}
\begin{note}
We have no log rules to split up \makeblank[1in]. If we have \makeblank[1in] (or \makeblank[1in]) inside a log, we have to leave it alone.
\end{note}
\begin{example}
Rewrite $\log_7\left(\dfrac{2x^4\sqrt{x-4}}{y^5z}\right)$ using the simplest possible logs.
\begin{lpic}{}{9}
\foreach \y in {-3,-5,-7,-9} \regtextleft{0}{\y}{${}={}$};
\end{lpic}
\end{example}
We notice the following, which gives us a shortcut.
\begin{itemize}
\item Factors in the \makeblank[1.5in] give us positive terms
\item Factors in the \makeblank[1.5in] give negative terms
\end{itemize}
Let's see another example, using that shortcut.
\begin{example}
Rewrite $\log_4\left(\dfrac{3(x-1)^5\sqrt[3]{y}}{z^4\sqrt{y+1}}\right)$ using the simplest possible logs.
\begin{lpic}{}{3}
\foreach \y in {-1,-3} \regtextleft{0}{\y}{${}={}$};
\end{lpic}
\end{example} 
\subsection{Log Rules to Combine} 
We will make extensive use of log rules to combine when solving equations.
\begin{example}
Use log rules to rewrite this expression as a single logarithm.
\[
2\log_5 x-3\log_5(y+7)-\tfrac{1}{3} \log_5 z + 4\log_5(x+1)
\]
We follow the same pattern as before:
\begin{itemize}
\item Negative terms go in the \makeblank[1.5in]
\item Positive terms go in the \makeblank[1.5in]
\end{itemize}
\begin{lpic}{}{3}
\foreach \y in {-1,-3} \regtextleft{0}{\y}{${}={}$};
\end{lpic}
\end{example}
We can only do this if all the logarithms have the same \makeblank[1in].
\begin{example}
Use log rules to combine these logarithms as much as possible.
\[
3\log_3 x-5\log_2 y+4\log_3 z+3\log_2 x
\]
\begin{lpic}{}{5}
\foreach \y in {-1,-3,-5} \regtextleft{0}{\y}{${}={}$};
\end{lpic}
\end{example} 
\subsection{Using the Log of a Power to Solve Equations} 
If we know that both sides of an equation are \makeblank[1in], we can take the log of both sides. This works especially well if one or both sides contains an \makeblank[1in] expression.
\begin{example}
Solve $3^x=5$.
\begin{itemize}
\item We cannot easily write 3 and 5 as \makeblank[1in] of the same number.
\item Therefore, we cannot solve this by \makeblank.
\item We'll take the \makeblank of both sides, so we can use a calculator.
\end{itemize}
\begin{lpic}{}{3}
\end{lpic}
\end{example}
We could also solve this by writing the equation in log form:
\[
\makeblank[3in]
\]
This suggests the last property of logarithms. 
\subsection{Change of Base for Logarithms} 
\begin{theorem}[Change of Base for Logarithms]
For any positive real numbers $b$, $c$, and $x$, with $b,c\neq 1$, we have:
\[
\log_b(x)=\frac{\log_c(x)}{\log_c(b)}
\]
\end{theorem}
\begin{proof}
Write $y=\log_b(x)$.
\begin{itemize}
\item Then $x=\makeblanksmall$
\item Take \makeblanksmall of both sides: $\makeblank[1in]=\makeblank[1in]$
\item Using the \makeblank rule, $\makeblank[1in]=\makeblank[1in]$
\item So $\makeblank[1in]=\makeblank[1in]$
\item Divide both sides by \makeblank[1in] to get $y=\dfrac{\blankparens[1in]}{\blankparens[1in]}$
\item Since both are equal to \makeblanksmall, $\log_b(x)=\frac{\log_c(x)}{\log_c(b)}$, as desired.
\end{itemize}
\end{proof}
At the moment, our primary use for this fact will be to calculate logs on the calculator. \emph{This is not its main purpose---you will need this property in calculus.}
\begin{Example}
Use the change of base rule to find $\log_7{212}$ on your calculator. Round your answer to three decimal places. Check your work by rewriting in exponential form.
\begin{lpic}{}{3}
\end{lpic}
\end{Example} 
%\subsection{Change of Base for Exponents} 
 
